\documentclass[10pt,a4paper]{article}
\usepackage{amsmath,amssymb,amsthm,graphicx}
% Mechanical presentation wrapper; no author mathematics added.
\usepackage{amsmath,amssymb,amsthm,bm,mathtools,hyperref}
\usepackage[margin=22mm]{geometry}
\newtheorem{theorem}{Theorem}
\newtheorem{proposition}[theorem]{Proposition}
\newtheorem{corollary}[theorem]{Corollary}
\theoremstyle{definition}
\newtheorem{definition}[theorem]{Definition}
\setlength{\emergencystretch}{3em}
\title{Polynomial-time reduction of max-average CSP to real semidefinite feasibility (Original Corollary 10)}
\author{Manuel Bodirsky, Georg Loho, Mateusz Skomra}
\date{doi:10.4230/LIPIcs.ICALP.2025.145; CC BY 4.0}
\begin{document}
\maketitle
\noindent\textit{Editorial scope.} One counted result: original Corollary 10. Theorem 6, Theorem 7, Proposition 8 and Theorem 9 are supporting author material, not additional counts. Original Sections 2–3 are retained to keep notation and the displayed system (1) readable. Corollary 11 is a direct consequence and is not counted.\par\medskip
\noindent\textit{Editorial note.} Source: Manuel Bodirsky, Georg Loho and Mateusz Skomra, Reducing Stochastic Games to Semidefinite Programming, ICALP 2025, LIPIcs 334, article 145, DOI 10.4230/LIPIcs.ICALP.2025.145. Licensed under CC BY 4.0, https://creativecommons.org/licenses/by/4.0/.\par\medskip
\noindent\textit{Editorial note.} This editable presentation mechanically transcribes the official published HTML and its mathematical alttext. Original numbering is preserved; LaTeXML nesting, footnote-marker duplication and the (P) reference are restored to the published PDF. No author mathematical statement or proof is added or corrected.\par\medskip
\noindent\textit{Editorial note.} The published sentence in Theorem 9 describing A\_i and B as k by k matrices is retained as written, alongside the original displayed block matrix and m=2(k+1). No mathematical correction is supplied.\par\medskip
\noindent\textit{Editorial scope.} Faithful collection and mechanical presentation of the fixed published author proof. No independent verification of author mathematics is claimed.\par\medskip
\noindent\textit{Allowed standard prerequisites (editorial).}\begin{itemize}
\item Real closedness of the Puiseux-series field. Original reference [11], Basu–Pollack–Roy, Section 2.6. Interpreting first-order formulas over Puiseux series and over the real field
\item Quantifier elimination and quantitative degree/coefficient bounds. Original reference [11], Theorems 2.77 and 14.6, and Corollary 10.4. Theorem 7 applies these standard bounds; the current author transfer proof is retained
\item Strong SDP duality under strict primal feasibility and a lower bound. Original Theorem 2; reference [12], Ben-Tal–Nemirovski, Theorem 2.4.1. Attained dual optimum in the current author singleton encoding of Theorem 9
\item PSD principal-minor test and block-diagonal conjunctions; elementary bit-size arithmetic. Original proof of Theorem 6 and Section 2.1; standard linear algebra and binary arithmetic. Transcribing polynomial-size LMI representations
\end{itemize}\medskip
% HTML id: S2.title; source line 727
\section*{Preliminaries}

% HTML id: S2.p1; source line 730
The set of natural numbers $\{0,1,2,\dots\}$ is denoted by ${\mathbb{N}}$, and the set of real numbers by ${\mathbb{R}}$. We assume familiarity with first-order logic; see, e.g., [38]. Let $\tau$ be a \emph{signature}, i.e., a set of function and relation symbols. A first-order $\tau$-formula is a formula built from the Boolean connectives, atomic $\tau$-formulas and quantification in the usual way. If $A$ is a $\tau$-structure and $\phi$ is a $\tau$-formula with free variables $x_{1},\dots,x_{k}$, then $\phi$ \emph{defines} over $A$ the relation


\[
\{(a_{1},\dots,a_{k})\in A^{k}\mid A\models\phi(a_{1},\dots,a_{k})\}.
\]
A set is called \emph{semialgebraic} if it has a first-order definition over the field $({\mathbb{R}};+,\cdot)$, allowing real parameters (we do not need the order in the signature since it is definable from addition and multiplication). A first-order formula $\phi(x_{1},\dots,x_{k})$ with free variables $x_{1},\dots,x_{k}$ is called \emph{primitive positive} if it is of the form


\[
\exists y_{1},\dots,y_{\ell}(\psi_{1}\wedge\cdots\wedge\psi_{m})
\]
where $\psi_{1},\dots,\psi_{m}$ are atomic formulas over the variables $x_{1},\dots,x_{k},y_{1},\dots,y_{\ell}$. A relation is called \emph{primitive positive definable in a $\tau$-structure $A$} if there exists a primitive positive $\tau$-formula which defines the relation over $A$.

% HTML id: S2.SS1.title; source line 760
\subsection*{Semidefinite Programming}

% HTML id: S2.SS1.p1; source line 763
For a symmetric matrix $S\in{\mathbb{R}}^{m\times m}$ we write $S\succeq 0$ if $S$ is \emph{positive semidefinite}, i.e., if $x^{\top}Sx\geq 0$ for all $x\in{\mathbb{R}}^{m}$, and $S\succ 0$ if it is \emph{positive definite}, i.e., if $x^{\top}Sx>0$ for all $x\in{\mathbb{R}}^{m}\setminus\{0\}$. A \emph{semidefinite program (SDP)} is an optimisation problem of the following form:


\begin{equation}\tag{P}\label{S2.Ex4}
\begin{aligned}
\displaystyle\inf & \displaystyle\sum_{i=1}^{n}c_{i}x_{i} \\
\text{such that} & \displaystyle\sum_{i=1}^{n}x_{i}A_{i}-B\succeq 0\,,
\end{aligned}
\end{equation}
where $A_{1},\dots,A_{n},B\in{\mathbb{R}}^{m\times m}$ are symmetric and $c_{1},\dots,c_{n}\in{\mathbb{R}}$. If the data $A_{1},\dots,A_{n},B$, $c_{1},\dots,c_{n}$ is rational, then the problem of deciding whether the infimum exists, and if so, to compute the infimum, is a computational problem known as the \emph{semidefinite programming problem}. The \emph{representation size} of the SDP then consists of the sum of the bit lengths of all the rational numbers in the input.

% HTML id: S2.SS1.p2; source line 787
The expression $\sum_{i=1}^{n}x_{i}A_{i}-B\succeq 0$ in \eqref{S2.Ex4} is called a \emph{linear matrix inequality (LMI)}, and defines the \emph{feasible region} of the SDP, which is called a \emph{spectrahedron}. The representation size of LMIs is defined analogously as for SDPs. We say that an SDP is \emph{strictly feasible} if $\sum_{i=1}^{n}x_{i}A_{i}-B\succ 0$ has a solution $x\in{\mathbb{R}}^{n}$.

% HTML id: Thmtheorem1; source line 793
\setcounter{theorem}{0}
\begin{definition}\label{Thmtheorem1}
Given symmetric matrices $A_{1},\dots,A_{n},B\in{\mathbb{Q}}^{m\times m}$ the \emph{feasibility problem for semidefinite programs} asks to decide whether there is a solution $x\in\mathbb{R}^{n}$ to an LMI


\[
x_{1}A_{1}+\dots+x_{n}A_{n}-B\succeq 0\enspace.
\]
\end{definition}

% HTML id: S2.SS1.p3; source line 809
Note that spectrahedra are convex and semialgebraic. It has been conjectured that every convex semialgebraic set is the projection of a spectrahedron (also called a \emph{spectrahedral shadow}), but this turned out to be false [55] (see [15] for an alternative proof of an even stronger inexpressibility result). It is easy to see that the set of spectrahedral shadows is closed under primitive positive definability, and that for every primitive positive formula $\phi$ we may find an LMI for the spectrahedral shadow defined by $\phi$ whose representation size is bounded by the sum of the representation sizes of conjuncts of $\phi$. We also note that the definition of a semidefinite program makes sense over any real closed field, not only over the real numbers. In particular, it is meaningful to consider semidefinite programs over non-Archimedean real closed fields, such as the field of Puiseux series which we discuss in Section 4.

% HTML id: S2.SS1.p4; source line 818
The \emph{(Frobenius) inner product} of two matrices $S,T\in{\mathbb{R}}^{m\times m}$ is defined as $\langle S,T\rangle:=\text{trace}(ST)=\sum_{i,j}S_{ij}T_{ij}$. The \emph{dual program} to the SDP from \eqref{S2.Ex4} (to which we refer as the \emph{primal}) is given by


\begin{equation}\tag{D}\label{S2.Ex6}
\begin{aligned}
\displaystyle\sup_{Y\succeq 0} & \displaystyle\langle B,Y\rangle \\
\text{such that} & \displaystyle\langle A_{i},Y\rangle=c_{i}\text{ for every }i\in\{1,\dots,n\}\,.
\end{aligned}
\end{equation}

% HTML id: S2.SS1.p5; source line 839
It is well-known that the value of the primal SDP is bounded from below by the value of the dual; however, the lower bound might not be tight, i.e., there might be a non-zero \emph{duality gap}. Moreover, both primal and dual might not have optimal solutions, i.e., the supremum may not be a maximum and the infimum might not be a minimum. See [54] for a more complicated, but exact dual. However, for us the following sufficient condition for a vanishing duality gap will be enough.

% HTML id: Thmtheorem2; source line 843
\setcounter{theorem}{1}
\begin{theorem}[{(see Theorem 2.4.1 in [ 12 ] ) }]\label{Thmtheorem2}
If the primal SDP is strictly feasible and bounded from below, then the dual program attains its supremum and there is no duality gap.

\end{theorem}

% HTML id: S2.SS2.title; source line 853
\subsection*{Stochastic Games}

% HTML id: S2.SS2.p1; source line 856
A \emph{stochastic mean payoff game} is a two-player game played on a finite edge-weighted directed graph $\vec{\mathcal{G}}=(V,E,r,p)$, i.e., $V$ is a finite set of nodes, $E\subseteq V^{2}$ is a set of edges, $r\colon E\to{\mathbb{Q}}$ is a \emph{weight function}, and $p\colon E\to{\mathbb{Q}}\cap[0,1]$ specifies \emph{edge probabilities}. The nodes of the graph are partitioned into three subsets, $V=V_{\mathrm{Min}}\sqcup V_{\mathrm{Max}}\sqcup V_{\mathrm{Rand}}$. The nodes in $V_{\mathrm{Min}}$ are controlled by player Min, the nodes in $V_{\mathrm{Max}}$ are controlled by player Max, and the nodes in $V_{\mathrm{Rand}}$ are controlled by nature; the probabilities of the outgoing edges of each vertex in $V_{\mathrm{Rand}}$ sum up to one. The players play the game by moving a pawn on the graph. When the pawn is on a node $v\in V_{\mathrm{Min}}$, player Min chooses an outgoing edge $(v,u)$ and moves the pawn to $u$. Analogously, player Max chooses the next move when the pawn is on a node $v\in V_{\mathrm{Max}}$. When the pawn is on a node $v\in V_{\mathrm{Rand}}$, the next move is chosen by nature randomly according to the probability distribution given by $p$. To ensure that players can always make a move, we suppose that all nodes of the graph have at least one outgoing edge. The weights of the edges of the graph represent payoffs that player Max receives from player Min after each move of the pawn. In order to state the objectives of both players, we first specify the strategies that they are allowed to use. A \emph{pure and positional} strategy (also called a \emph{policy}) for player Min is a function $\sigma\colon V_{\mathrm{Min}}\to V$ that to each node controlled by Min associates a neighbor, $(v,\sigma(v))\in E$. Player Min plays according to $\sigma$ if they always move the pawn to $\sigma(v)$ when the pawn lands on $v$. We analogously define a policy $\tau\colon V_{\mathrm{Max}}\to V$ of player Max. Note that if we fix a pair of policies and an initial position of the pawn, then the resulting movement of the pawn is a Markov chain on the graph $(V,E)$. Given an initial position $u\in V$ of the pawn and a pair of policies $(\sigma,\tau)$, the \emph{mean payoff reward} of player Max is the quantity


\[
g(\sigma,\tau)_{u}\coloneqq\lim_{N\to\infty}\frac{1}{N}\mathbb{E}^{u}_{\sigma,\tau}(r_{u_{0}u_{1}}+\dots+r_{u_{N-1}u_{N}})\,,
\]
where $u=u_{0},u_{1},\dots,u_{N}$ is the path of the pawn in the first $N$ turns of the game and the expectation is taken with respect to the Markov chain obtained by fixing $(\sigma,\tau)$. The ergodic theorem of Markov chains shows that $g(\sigma,\tau)_{u}$ is well-defined. The objective of player Max is to maximize the mean payoff reward while the objective of player Min is to minimize this reward. It is known that both players in a stochastic mean payoff game have optimal policies [46]. In other words, there exist policies $(\sigma^{*},\tau^{*})$ and a vector $\chi\in\mathbb{R}^{V}$ such that for every $u$ and every $(\sigma,\tau)$ we have the inequality


\[
g(\sigma^{*},\tau)_{u}\leq\chi_{u}\leq g(\sigma,\tau^{*})_{u}\,.
\]
The reward $\chi=g(\sigma^{*},\tau^{*})$ is called the \emph{value} of the game. In order to make these games computationally tractable, we suppose that the weights and the probability distributions at nodes in $V_{\mathrm{Rand}}$ are rational. The nodes of the graph are also called the \emph{states} of the game.

% HTML id: S2.SS2.p2; source line 885
A \emph{simple stochastic game} is a special case of a stochastic mean payoff game. In this game, every node in $V_{\mathrm{Rand}}$ has out-degree exactly two and the probability distribution at every such node is $(1/2,1/2)$. Furthermore, the graph is equipped with two special terminal states called $\mathrm{Win}$ and $\mathrm{Lose}$. These states are absorbing, i.e., $\mathrm{Win}$ has only one outgoing edge that goes back to $\mathrm{Win}$ and $\mathrm{Lose}$ has only one outgoing edge that goes back to $\mathrm{Lose}$. The weight of the edge going from $\mathrm{Win}$ to $\mathrm{Win}$ is equal to $1$, while the weight of all of the other edges of the graph is equal to $0$. In this case, the mean payoff reward of player Max is simply the probability that the pawn reaches the winning state, $g(\sigma,\tau)_{u}=\mathbb{P}^{u}_{\sigma,\tau}(\text{Reach}(\mathrm{Win}))$.

% HTML id: S2.SS2.p3; source line 888
A simple stochastic game is called \emph{stopping} if for every initial position $u\in V$ and every pair of policies $(\sigma,\tau)$, the pawn reaches a terminal state with probability one, $\mathbb{P}^{u}_{\sigma,\tau}(\text{Reach}(\mathrm{Win},\mathrm{Lose}))=1$. In other words, the game is stopping if the Markov chain obtained by fixing any pair of policies has only two recurrent classes, namely $\{\mathrm{Win}\}$ and $\{\mathrm{Lose}\}$.

% HTML id: S2.SS2.p4; source line 891
The computational problem associated with stochastic mean payoff games is the following: given a stochastic mean payoff game, find its value and a pair of optimal policies. In [8], Andersson and Miltersen proved that this problem is polynomial-time (Turing) reducible to the problem of finding the value of a stopping simple stochastic game. By a binary search argument, the problem of computing the value of a stopping simple stochastic game can be further reduced to the decision problem, in which we want to decide whether the value of a given state of such a game is at least $1/2$. The details of this reduction can be found, e.g., in [6] or [57, Section 7.1].

% HTML id: Thmtheorem3; source line 896
\setcounter{theorem}{2}
\begin{theorem}[{( [ 6 , Lemma 48] ) }]\label{Thmtheorem3}
The problem of computing the value and a pair of optimal policies of a stochastic mean payoff game is polynomial-time (Turing) reducible to the problem of deciding whether a given state of a stopping simple stochastic game has value at least $1/2$.

\end{theorem}

% HTML id: S3.title; source line 907
\section*{From stochastic games to max-average constraints}

% HTML id: S3.p1; source line 911
In this section, we focus on the decision problem stated in Theorem 3 above: given a stopping simple stochastic game we want to decide whether the value of a given state is at least $1/2$. This problem can be reduced to max-average constraint satisfaction problems (\emph{max-average CSPs}). The ideas for this reduction are not new, but we are not aware of an explicit reference for this fact, and we therefore present a short proof in this section. In the next section, we will show that max-average CSPs then further reduce to the feasibility problem of non-Archimedean semidefinite programs. The fact that stochastic games can be encoded as non-Archimedean SDPs was proven in [6] and our reduction is very similar. While their approach relies on reducing stochastic games to ergodic games, we omit this step and obtain a slightly more straightforward reduction.

% HTML id: Thmtheorem4; source line 920
\setcounter{theorem}{3}
\begin{definition}\label{Thmtheorem4}
An instance of a \emph{max-average CSP} is a conjunction of constraints, each of which has one of the following forms:


\begin{enumerate}
\item $x_{0}\leq\max(x_{1},\dots,x_{k})$ where $k\in\mathbb{N}$,


\item $x_{0}\leq\frac{x_{1}+x_{2}}{2}$,


\item $x_{0}=c$ where $c$ is a rational constant.


\end{enumerate}
The computational task is to determine whether such a conjunction has a solution over $\mathbb{Q}\cup\{-\infty\}$ (equivalently, over $\mathbb{R}\cup\{-\infty\}$).

\end{definition}

% HTML id: S3.p2; source line 949
Note that max-average CSP instances may also use constraints of the form $x_{0}\leq\min(x_{1},\dots,x_{k})$ where $k\in\mathbb{N}$, $x_{0}\leq c$, or $x_{0}\geq c$ for some rational constant $c$, because the respective relations are primitively positively definable over the other relations.

% HTML id: S3.p3; source line 952
We now show that the problem of solving stochastic mean payoff games can be reduced to max-average CSPs. A reduction to slightly more general CSPs, including constraints of the form $x_{0}\leq c+x_{1}$, follows from [6, Theorem 18] and is based on a Collatz–Wielandt property of order-preserving and additively homogeneous maps. However, for our purposes, it is useful that the reduction does not use these additional constraints. Instead, we give a proposition that characterizes the value of a stopping simple stochastic game. It was already proven in [22] that this value is a unique fixed point of an operator involving min, max, and averaging operations. By replacing the fixed point equations with inequalities, we obtain a set of certificates which prove that the value is at least $1/2$. This observation is proven, e.g., in [13] and gives inequalities very similar to the inequalities on Shapley operators used in [6].

% HTML id: Thmtheorem5; source line 959
\setcounter{theorem}{4}
\begin{proposition}\label{Thmtheorem5}
Consider a stopping simple stochastic game played on a directed graph $\vec{\mathcal{G}}=([n],E)$. Then, the value of a state $k\in[n]$ is at least $1/2$ if and only if the following system of inequalities (an instance of a max-average constraint satisfaction problem) has a solution $x\in\mathbb{Q}_{\geq 0}^{n}$:


\begin{equation}\tag{1}\label{S3.E1}
\begin{aligned}
\displaystyle x_{i} & \displaystyle\leq\max_{(i,j)\in E}\{x_{j}\} &  & \displaystyle\text{for every }i\in V_{\mathrm{Max}}, \\
\displaystyle x_{i} & \displaystyle\leq\min_{(i,j)\in E}\{x_{j}\} &  & \displaystyle\text{for every }i\in V_{\mathrm{Min}}, \\
\displaystyle x_{i} & \displaystyle\leq\frac{1}{2}\sum_{(i,j)\in E}x_{j} &  & \displaystyle\text{for every }i\in V_{\mathrm{Rand}}, \\
\displaystyle x_{\mathrm{Lose}} & \displaystyle=0, \\
\displaystyle x_{\mathrm{Win}} & \displaystyle=1, \\
\displaystyle x_{k} & \displaystyle\geq 1/2.
\end{aligned}
\end{equation}
\end{proposition}

% HTML id: S3.SS0.SSS0.P0.SPx1; source line 1014
\begin{proof}
Let $(Q)$ denote the system of inequalities obtained from \eqref{S3.E1} by removing the inequality $x_{k}\geq 1/2$. By [22] or [13, Theorem 77], the value of the game satisfies $(Q)$ with inequalities replaced by equalities. Hence, if the value at state $k$ is at least $1/2$, then the value of the game satisfies \eqref{S3.E1}. Moreover, the value of the game is a rational vector by [22, Lemma 2] or [13, Corollary 18]. Conversely, if a point $x\in\mathbb{Q}_{\geq 0}^{n}$ satisfies $(Q)$, then [13, Lemma 83] implies that for all $i\in[n]$, the entry $x_{i}$ is not greater than the value of the game at state $i$. Hence, if we further suppose that $x_{k}\geq 1/2$, then the value of the game at state $k$ is at least $1/2$. 

\end{proof}

% HTML id: S4.title; source line 1024
\section*{From max-average constraints to non-Archimedean SDPs}

% HTML id: S4.p1; source line 1027
By “lifting” the inequalities \eqref{S3.E1} to Puiseux series, we get the reduction from max-average constraint satisfaction problems to non-Archimedean SDPs. A \emph{(formal real) Puiseux series} is a series of the form


\[
\bm{x}=\bm{x}(t)=c_{0}t^{a/n}+c_{1}t^{(a-1)/n}+c_{2}t^{(a-2)/n}+\dots\,,
\]
where $t$ is a formal parameter, the number $n$ is a positive integer, $a$ is an integer, and the coefficients $c_{i}$ are real. We further assume that $c_{0}\neq 0$ for every Puiseux series except the zero series. The number $a/n$ is called the \emph{valuation} of the series $\bm{x}$ and is denoted by $\operatorname*{\mathsf{val}}(\bm{x})$, with the convention that $\operatorname*{\mathsf{val}}(0)=-\infty$. Puiseux series can be added and multiplied in the natural way. Moreover, they can be ordered by defining $\bm{x}\geq 0$ if $c_{0}\geq 0$ and $\bm{x}\geq\bm{y}$ if $\bm{x}-\bm{y}\geq 0$. It is known that Puiseux series form a real closed field, see, e.g., [11, Section 2.6] for a detailed discussion.\footnote{In this work, we use a convention in which one interprets $t$ as a “very large” parameter. Many authors use the opposite but equivalent convention in which $t$ is a very small parameter. This explains the differences in notation between our work and [11].} This field is non-Archimedean, meaning that the order does not satisfy the Archimedean property: the series $\bm{x}(t)=t$ is larger than any natural number, i.e., $t\geq m$ for all $m\in\mathbb{N}$. We denote the field of Puiseux series by $\mathbb{K}$. For any $\bm{x},\bm{y}\in\mathbb{K}$, the following properties of the valuation function are immediate from the definitions:


\begin{equation}\tag{2}\label{S4.E2}
\begin{aligned}
\displaystyle\operatorname*{\mathsf{val}}(\bm{x}+\bm{y}) & \displaystyle\leq\max\{\operatorname*{\mathsf{val}}(\bm{x}),\operatorname*{\mathsf{val}}(\bm{y})\}\,, \\
\displaystyle\operatorname*{\mathsf{val}}(\bm{x}\bm{y}) & \displaystyle=\operatorname*{\mathsf{val}}(\bm{x})+\operatorname*{\mathsf{val}}(\bm{y})\,, \\
\displaystyle\bm{x}\geq\bm{y}\geq 0 & \displaystyle\implies\operatorname*{\mathsf{val}}(\bm{x})\geq\operatorname*{\mathsf{val}}(\bm{y})\,.
\end{aligned}
\end{equation}

% HTML id: Thmtheorem6; source line 1066
\setcounter{theorem}{5}
\begin{theorem}\label{Thmtheorem6}
An instance $\phi$ with variables $x_{1},\dots,x_{n}$ of a max-average constraint satisfaction problem has a solution in $(\mathbb{Q}\cup\{-\infty\})^{n}$ if and only if the following system of inequalities $S_{\phi}(t)$, which is a (quadratic) semidefinite program over Puiseux series, has a solution in $\mathbb{K}^{n}$: for every variable $x_{i}$ of $\phi$ we have a variable $\bm{x}_{i}$, and the inequalities are


\begin{equation}\tag{3}\label{S4.E3}
\begin{aligned}
\displaystyle\bm{x}_{i_{0}} & \displaystyle\leq\bm{x}_{i_{1}}+\cdots+\bm{x}_{i_{k}} &  & \text{for every constraint }x_{i_{0}}\leq\max(x_{i_{1}},\dots,x_{i_{k}})\text{ in }\phi\text{,} \\
\displaystyle\bm{x}^{2}_{i_{0}} & \displaystyle\leq\bm{x}_{i_{1}}\bm{x}_{i_{2}} &  & \text{for every constraint }x_{i_{0}}\leq\frac{x_{i_{1}}+x_{i_{2}}}{2}\text{ in }\phi\text{,} \\
\displaystyle\bm{x}_{i_{0}} & \displaystyle=t^{c} &  & \text{for every constraint }x_{i_{0}}=c\text{ in }\phi\text{,} \\
\displaystyle\bm{x}_{i_{0}} & \displaystyle\geq 0 &  & \text{for every }i_{0}\in\{1,\dots,n\}\text{.}
\end{aligned}
\end{equation}
\end{theorem}

% HTML id: S4.SS0.SSS0.P0.SPx1; source line 1121
\begin{proof}
Suppose that $\bm{x}=(\bm{x}_{1},\dots,\bm{x}_{n})\in\mathbb{K}^{n}$ satisfies \eqref{S4.E3}. Then, $\operatorname*{\mathsf{val}}(\bm{x})$ belongs to $(\mathbb{Q}\cup\{-\infty\})^{n}$ and satisfies $\phi$ by the properties of the valuation function \eqref{S4.E2}. Conversely, if $(x_{1},\dots,x_{n})\in(\mathbb{Q}\cup\{-\infty\})^{n}$ satisfies $\phi$, then $(t^{x_{1}},\dots,t^{x_{n}})$ (with the convention that $t^{-\infty}=0$) satisfies \eqref{S4.E3}.

This amounts to the feasibility problem for a semidefinite program: the linear (in)equalities can be expressed by diagonal matrices and each quadratic condition in combination with the non-negativity of all variables is just the positive semidefiniteness of a symmetric 2-by-2 matrix. Indeed, the symmetric matrix $\begin{pmatrix}x_{1}&x_{3}\\ x_{3}&x_{2}\end{pmatrix}$ is positive semidefinite if and only if $x_{1}\geq 0$, $x_{2}\geq 0$, and $x_{1}x_{2}\geq x_{3}^{2}$. (More generally, a symmetric matrix is positive semidefinite if and only if all of its principal minors are nonnegative.) These are precisely the inequalities that appear in \eqref{S4.E3}, together with some linear inequalities. A conjunction of such inequalities can then be expressed by a single LMI given by block-diagonal matrices (a symmetric block-diagonal matrix is positive semidefinite if and only if every block is positive semidefinite). 

\end{proof}

% HTML id: S5.title; source line 1140
\section*{From non-Archimedean to real SDPs}

% HTML id: S5.p1; source line 1143
To conclude, we aim to reduce the non-Archimedean semidefinite program from \eqref{S4.E3} to a real semidefinite program. We will do this in two steps. First, we replace the formal parameter $t$ by a real number that is sufficiently large (doubly exponential in the size of the instance). In fact, any non-Archimedean semialgebraic feasiblity problem can be reduced to a semialgebraic feasibility problem over the reals by replacing $t$ with a sufficiently large real number. This follows from the quantifier elimination in real closed fields coupled with the definition of the order in Puiseux series. By using the effective bounds on quantifier elimination discussed, e.g., in [11], we can give a bound on this “sufficiently large” number. Second, we will show that this SDP with very large coefficients can be reduced to another SDP that has coefficients of polynomial bit-size.

% HTML id: S5.p2; source line 1148
Before presenting the proofs, we note that a reduction from non-Archimedean SDPs coming from games to real SDPs with large coefficients was already studied in [6, Theorem 36]. However, the bounds given in [6] rely on an additional assumption on the underlying game. In particular, they cannot be applied directly to the SDP given in \eqref{S4.E3}. For max-average systems coming from games, one could overcome this restriction by perturbing the underlying game slightly. Nevertheless, the bounds coming from quantifier elimination which we present below do not require any additional assumptions and can be of independent interest. We recall that Tarski’s quantifier elimination allows to eliminate quantifiers of first-order $\tau$-formulas, with signature $\tau$ given by $\tau=\{+,\cdot,0,1,\leq\}$.

% HTML id: Thmtheorem7; source line 1151
\setcounter{theorem}{6}
\begin{theorem}\label{Thmtheorem7}
Let $P_{1},\dots,P_{m}\in\mathbb{Z}[x_{0},x_{1},\dots,x_{n}]$ be a collection of polynomials with integer coefficients and let $\diamond\in\{\leq,\geq,<,>,=,\neq\}^{m}$ be a sign pattern. Consider the semialgebraic set over Puiseux series defined by


\[
\bm{S}\coloneqq\{\bm{x}\in\mathbb{K}^{n}\colon\forall i\in[m],\,P_{i}(t,\bm{x})\diamond_{i}0\}\,.
\]
Furthermore, for any real $L>0$ consider the semialgebraic set over the real numbers defined by


\[
S(L)\coloneqq\{x\in\mathbb{R}^{n}\colon\forall i\in[m],\,P_{i}(L,x)\diamond_{i}0\}\,.
\]
Then, $\bm{S}$ is nonempty if and only if $S(L)$ is nonempty for all sufficiently large $L$. More precisely, there exists an absolute constant $C>0$ with the following property. If the maximal degree of any polynomial $P_{i}$ is bounded by $d\geq 2$ and $s$ denotes the maximal bit-size of any coefficient of $P_{i}$, then either for all $L>2^{sd^{Cn}}$ we have $S(L)=\emptyset$, or for all $L>2^{sd^{Cn}}$ we have $S(L)\neq\emptyset$. In the former case $\bm{S}$ is empty and in the latter case $\bm{S}$ is nonempty.

\end{theorem}

% HTML id: S5.SS0.SSS0.P0.SPx1; source line 1184
\begin{proof}
Both sets $S(L)$ and $\bm{S}$ can be defined by the same first-order formula with a free variable $x_{0}$ that can be substituted by a real number $L$ (to obtain $S(L)$) or a Puiseux series $t$ (to obtain $\bm{S}$). Let $\tau=\{+,\cdot,0,1,\leq\}$ and consider the first-order $\tau$-formula $\phi(x_{0})$ given by


\[
\exists x_{1},\dots,\exists x_{n}\operatorname*{\bigwedge}_{i=1}^{m}(P_{i}(x_{0},x)\diamond_{i}0)\,.
\]
By Tarski’s quantifier elimination theorem [11, Theorem 2.77], there exists a quantifier-free formula $\psi(x_{0})$ that is equivalent to $\phi(x_{0})$ in the theory of real closed fields. The bounds on quantifier elimination given in [11, Theorem 14.6] show that $\psi(x_{0})$ can be taken to be of the form


\[
\operatorname*{\bigvee}_{i=1}^{I}\operatorname*{\bigwedge}_{j=1}^{J_{i}}\left(\operatorname*{\bigvee}_{k=1}^{N_{ij}}\operatorname*{\mathrm{sign}}(Q_{ijk}(x_{0}))=\delta_{ijk}\right)\,,
\]
where $\delta_{ijk}\in\{-1,0,1\}$ and $Q_{ijk}\in\mathbb{Z}[x_{0}]$ are univariate polynomials with integer coefficients. Furthermore, every $Q_{ijk}$ has degree bounded by $d^{O(n)}$ and every coefficient of every $Q_{ijk}$ has bit-size bounded by $sd^{O(n)}$. Hence, the Cauchy bound [11, Corollary 10.4] implies that every real root of $Q_{ijk}$ has absolute value bounded by $2^{sd^{O(n)}}$ (provided that $Q_{ijk}$ is a non-zero polynomial). In particular, the sign of $Q_{ijk}(L)$ is the same for all $L>2^{sd^{O(n)}}$. Therefore, either $S(L)$ is empty for all such $L$ or $S(L)$ is nonempty for all such $L$. Moreover, $Q_{ijk}(t)\in\mathbb{K}$ is a Puiseux series with finitely many terms. By definition, the sign of this series is the same as the sign of $Q_{ijk}(L)\in\mathbb{R}$. Hence, $\bm{S}$ is nonempty if and only if $S(L)$ is nonempty. 

\end{proof}
In the next proposition and later, we use the following notation. Fix an instance $\phi$ of a max-average CSP and let $(c_{i})_{i}$ be the rational constants used in the constraints of $\phi$. For every $i$ we denote $c_{i}=p_{i}/q_{i}$ where $p_{i}\in\mathbb{Z}$ and $q_{i}\in\mathbb{N}^{*}$ and we put $W\coloneqq 2+\max_{i}\{|p_{i}|+q_{i}\}$.

% HTML id: Thmtheorem8; source line 1218
\setcounter{theorem}{7}
\begin{proposition}\label{Thmtheorem8}
There exists an absolute integer constant $M>0$ such that an instance $\phi$ of a max-average constraint satisfaction problem has a solution if and only if the real SDP $S_{\phi}(L)$ obtained from $S_{\phi}(t)$ (see Theorem 6) by replacing $t$ with $L>2^{2^{Mn\log_{2}W}}$, is feasible.

\end{proposition}

% HTML id: S5.SS0.SSS0.P0.SPx2; source line 1230
\begin{proof}
Consider the system of inequalities given in \eqref{S4.E3}. Since $\bm{x}_{i}\geq 0$ for all $i$, any equality of the form $\bm{x}_{i}=t^{c_{j}}=t^{p_{j}/q_{j}}$ can be replaced with an equivalent equality $\bm{x}_{i}^{q_{j}}=t^{p_{j}}$ if $p_{j}\geq 0$ and $t^{-p_{j}}\bm{x}_{i}^{q_{j}}=1$ if $p_{j}<0$. If we further replace the parameter $t$ by a variable symbol $x_{0}$, \eqref{S4.E3} becomes a system of polynomial inequalities of degree at most $W$ and such that the coefficients of the polynomials appearing in these inequalities are integers in $\{-1,0,1\}$. Hence, by Theorem 7, there exists $M>0$ such that the system \eqref{S4.E3} has a solution over Puiseux series if and only if the system obtained by replacing $t$ with $L>2^{2^{Mn\log_{2}W}}$ has a solution over the real numbers. Therefore, the claim follows from Theorem 6. 

\end{proof}
In Proposition 8, we obtain a (real) semidefinite program, whose size, however, is exponential in the size of $\phi$ and therefore does not provide a polynomial-time reduction to the semidefinite program feasiblity problem. However, by the following result we can also reduce to a semidefinite program of polynomial representation size.

% HTML id: Thmtheorem9; source line 1241
\setcounter{theorem}{8}
\begin{theorem}\label{Thmtheorem9}
For every $n\in{\mathbb{Z}}$, the singleton set $\{2^{n}\}$ is the projection of a spectrahedron whose representation size is polynomial in the bit-size of $n$.

\end{theorem}

% HTML id: S5.SS0.SSS0.P0.SPx3; source line 1250
\begin{proof}
To prove the claim, we extend the constructions of SDPs with solutions of doubly-exponential size given in [51, 54]. For $n=0$, the set $\{x\in\mathbb{R}\mid x=1\}$ is the desired spectrahedron. Suppose that $n\neq 0$ and let $|n|=\sum_{i=0}^{k}\bar{b}_{k-i}2^{i}$ be the binary representation of $|n|$ for $\bar{b}_{i}\in\{0,1\}$. Let $b_{i}=\operatorname*{\mathrm{sign}}(n)\bar{b}_{i}\in\{-1,0,1\}$ for all $i$, so that $n=\sum_{i=0}^{k}b_{k-i}2^{i}$. We first show that the set $\{x\in{\mathbb{R}}\mid x\geq 2^{n}\}$ is the projection of a spectrahedron $S$ of representation size that is polynomial in the bit-size of $n$. Let $S$ be defined by the linear matrix inequality


\[
\displaystyle A_{0}x_{0}+\cdots+A_{k}x_{k}-B=\begin{pmatrix}x_{0}&1&\\ 1&2^{-b_{0}}\\ &&x_{1}&x_{0}\\ &&x_{0}&2^{-b_{1}}\\ &&&&\ddots\\ &&&&&x_{k}&x_{k-1}\\ &&&&&x_{k-1}&2^{-b_{k}}\end{pmatrix}\succeq 0\,.
\]
This LMI expresses that


\[
x_{0}\geq 2^{b_{0}},\,x_{1}\geq 2^{b_{1}}x_{0}^{2},\,\dots,\,x_{k}\geq 2^{b_{k}}x_{k-1}^{2}\,.
\]
By induction, for every $\ell\geq 0$ we get $x_{\ell}\geq 2^{b_{0}2^{\ell}+b_{1}2^{\ell-1}+\dots+b_{\ell}2^{0}}$. Moreover, the vector $(x_{0},\dots,x_{k})$ with $x_{\ell}=2^{b_{0}2^{\ell}+b_{1}2^{\ell-1}+\dots+b_{\ell}2^{0}}$ for every $\ell\in\{0,\dots,k\}$ satisfies this LMI, so we obtain the desired set by projecting $S$ to the last coordinate. We use strong duality to obtain the singleton set $\{2^{n}\}$ as follows. Consider the SDP


\[
\inf\,(0\cdot x_{0}+\cdots+0\cdot x_{k-1}+1\cdot x_{k})
\]
over the spectrahedron $S$ specified above. Then, the dual is of the form $\sup\,\langle B,Y\rangle$ such that $\langle A_{i},Y\rangle=c_{i}$ for every $i\in\{0,\dots,k\}$ and $Y\succeq 0$, where $(c_{0},\dots,c_{k})=(0,\dots,0,1)$. Clearly, the primal is strictly feasible and bounded by $2^{n}$ from below, so we may use strong duality (Theorem 2) and obtain that $\sup\,\langle B,Y\rangle=2^{n}$ and that the supremum is attained for some $Y$. Then the following primitive positive formula over spectrahedral sets with one free variable $x_{k}$ defines the singleton relation $\{2^{n}\}$:


\begin{equation}\tag{4}\label{S5.E4}
\begin{aligned}
\displaystyle\exists x_{0},\dots,x_{k-1},y_{1,1},\dots,y_{m,m}:A_{0}x_{0}+\cdots+A_{k}x_{k}-B\succeq 0 &  \\
\displaystyle\wedge\,\langle A_{0},Y\rangle=0\,\wedge\cdots\wedge\,\langle A_{k-1},Y\rangle=0\,\wedge\,\langle A_{k},Y\rangle=1 & \displaystyle\wedge Y\succeq 0 \\
 & \displaystyle\wedge\langle B,Y\rangle=x_{k}
\end{aligned}
\end{equation}
where $Y$ denotes $(y_{i,j})_{i,j\in\{1,\dots,m\}}$ for $m=2(k+1)$. As we have mentioned before, spectrahedral shadows are closed under primitive positive definability, and the representation size of the resulting LMI is polynomially bounded by the representation sizes for the conjuncts. Indeed, the matrices $A_{i}$ and $B$ have coefficients in ${0,\pm 1/2,\pm 1,\pm 2}$ and are of size $k\times k$ which is polynomial in the bit-size of $n$. Hence, the system of matrix inequalities \eqref{S5.E4} can be encoded using a number of bits that is polynomial in the bit-size of $n$. This system of inequalities can be then encoded by a single block-diagonal LMI as in the proof of Theorem 6. This increases the number of necessary bits by a polynomial factor. We thus obtain an LMI of representation size that is polynomial in $k$. 

\end{proof}

% HTML id: Thmtheorem10; source line 1330
\setcounter{theorem}{9}
\begin{corollary}\label{Thmtheorem10}
There is a polynomial-time reduction from the max-average constraint satisfaction problem to the feasibility problem for semidefinite programs.

\end{corollary}

% HTML id: S5.SS0.SSS0.P0.SPx4; source line 1339
\begin{proof}
We have already specified a first reduction to a (real) SDP in Proposition 8. To obtain a proper polynomial-time reduction, we modify this SDP as follows. Fix an instance $\phi$; let $c_{i}$, $p_{i}$, $q_{i}$, and $W$ be as introduced before Proposition 8. Define $D\coloneqq\prod_{i}q_{i}$ and $K\coloneqq 2DW^{Mn}=D2^{1+Mn\log_{2}W}$, where $M$ is the integer constant from Proposition 8. Furthermore, let $L\coloneqq 2^{K}$. By Proposition 8, $\phi$ has a solution if and only if $S_{\phi}(L)$ has a solution over the real numbers. This SDP has constraints of the form $x_{i}=L^{c_{j}}=2^{c_{j}K}$. Since $K$ is an integer divisible by $D$, the number $c_{j}K$ is an integer for all $j$. Hence, by introducing new variables, we may replace the constraints $x_{j}=2^{c_{j}K}$ by the corresponding LMIs of small representation size from Theorem 9. The resulting SDP can be computed in polynomial time. Furthermore, it is feasible if and only if the original SDP is feasible, which proves the statement. 

\end{proof}

\begin{thebibliography}{99}
\bibitem[1]{bib1} M. Akian, S. Gaubert, and A. Guterman. Tropical polyhedra are equivalent to mean payoff games. Int. J. Algebra Comput., 22(1):125001 (43 pages), 2012. \href{https://doi.org/10.1142/S0218196711006674}{doi:10.1142/S0218196711006674}.

\bibitem[2]{bib2} Marianne Akian, Ravindra Bapat, and Stéphane Gaubert. Max-plus algebra. In Chapter 25 in “Handbook of linear algebra”, volume 39 of Discrete Mathematics and Its Applications. Chapman and Hall, 2007.

\bibitem[3]{bib3} X. Allamigeon, P. Benchimol, S. Gaubert, and M. Joswig. Combinatorial simplex algorithms can solve mean payoff games. SIAM J. Optim., 24(4):2096–2117, 2014. \href{https://doi.org/10.1137/140953800}{doi:10.1137/140953800}.

\bibitem[4]{bib4} X. Allamigeon, P. Benchimol, S. Gaubert, and M. Joswig. Tropicalizing the simplex algorithm. SIAM J. Discrete Math., 29(2):751–795, 2015. \href{https://doi.org/10.1137/130936464}{doi:10.1137/130936464}.

\bibitem[5]{bib5} X. Allamigeon, P. Benchimol, S. Gaubert, and M. Joswig. What tropical geometry tells us about the complexity of linear programming. SIAM Rev., 63(1):123–164, 2021. \href{https://doi.org/10.1137/20m1380211}{doi:10.1137/20m1380211}.

\bibitem[6]{bib6} X. Allamigeon, S. Gaubert, and M. Skomra. Solving generic nonarchimedean semidefinite programs using stochastic game algorithms. J. Symbolic Comput., 85:25–54, 2018. \href{https://doi.org/10.1016/j.jsc.2017.07.002}{doi:10.1016/j.jsc.2017.07.002}.

\bibitem[7]{bib7} Eric Allender, Peter Bürgisser, Johan Kjeldgaard-Pedersen, and Peter Bro Miltersen. On the complexity of numerical analysis. SIAM J. Comput., 38(5):1987–2006, 2008/09. \href{https://doi.org/10.1137/070697926}{doi:10.1137/070697926}.

\bibitem[8]{bib8} D. Andersson and P. B. Miltersen. The complexity of solving stochastic games on graphs. In Proceedings of the 20th International Symposium on Algorithms and Computation (ISAAC), volume 5878 of Lecture Notes in Comput. Sci., pages 112–121. Springer, 2009. \href{https://doi.org/10.1007/978-3-642-10631-6_13}{doi:10.1007/978-3-642-10631-6\_13}.

\bibitem[9]{bib9} A. Atserias and E. Maneva. Mean-payoff games and the max-atom problem, 2009. URL: \href{https://www.cs.upc.edu/~atserias/papers/mean-payoff-games-max-atom/mpma.pdf}{https://www.cs.upc.edu/˜atserias/papers/mean-payoff-games-max-atom/mpma.pdf}.

\bibitem[10]{bib10} David Avis and Oliver Friedmann. An exponential lower bound for Cunningham’s rule. Math. Program., 161(1-2):271–305, 2017. \href{https://doi.org/10.1007/S10107-016-1008-4}{doi:10.1007/S10107-016-1008-4}.

\bibitem[11]{bib11} S. Basu, R. Pollack, and M.-F. Roy. Algorithms in Real Algebraic Geometry, volume 10 of Algorithms Comput. Math. Springer, Berlin, 2006. \href{https://doi.org/10.1007/3-540-33099-2}{doi:10.1007/3-540-33099-2}.

\bibitem[12]{bib12} A. Ben-Tal and A. Nemirovski. Lectures on Modern Convex Optimization: Analysis, Algorithms, and Engineering Applications, volume 2 of MOS-SIAM Ser. Optim. SIAM, Philadelphia, PA, 2001. \href{https://doi.org/10.1137/1.9780898718829}{doi:10.1137/1.9780898718829}.

\bibitem[13]{bib13} N. Bertrand, P. Bouyer-Decitre, N. Fijalkow, and M. Skomra. Stochastic games. In N. Fijalkow, editor, Games on Graphs. 2023.

\bibitem[14]{bib14} Marc Bezem, Robert Nieuwenhuis, and Enric Rodríguez-Carbonell. The max-atom problem and its relevance. In Proceedings of the International Conference on Logic for Programming, Artificial Intelligence, and Reasoning (LPAR), pages 47–61, 2008. \href{https://doi.org/10.1007/978-3-540-89439-1_4}{doi:10.1007/978-3-540-89439-1\_4}.

\bibitem[15]{bib15} Manuel Bodirsky, Mario Kummer, and Andreas Thom. Spectrahedral shadows and completely positive maps on real closed fields. J. Eur. Math. Soc., 2024.

\bibitem[16]{bib16} Manuel Bodirsky and Marcello Mamino. Tropically convex constraint satisfaction. Theory of Computing Systems, 62(3):481–509, 2018. An extended abstract of the paper appeared under the title “Max-Closed Semilinear Constraints” in the proceedings of CSR’16; preprint available under ArXiv:1506.04184.

\bibitem[17]{bib17} Endre Boros, Khaled Elbassioni, Vladimir Gurvich, and Kazuhisa Makino. A convex programming-based algorithm for mean payoff stochastic games with perfect information. Optim. Lett., 11(8):1499–1512, 2017. \href{https://doi.org/10.1007/s11590-017-1140-y}{doi:10.1007/s11590-017-1140-y}.

\bibitem[18]{bib18} Julian C. Bradfield and Igor Walukiewicz. The mu-calculus and model checking. In Edmund M. Clarke, Thomas A. Henzinger, Helmut Veith, and Roderick Bloem, editors, Handbook of Model Checking, pages 871–919. Springer, 2018. \href{https://doi.org/10.1007/978-3-319-10575-8_26}{doi:10.1007/978-3-319-10575-8\_26}.

\bibitem[19]{bib19} Cristian S. Calude, Sanjay Jain, Bakhadyr Khoussainov, Wei Li, and Frank Stephan. Deciding parity games in quasi-polynomial time. SIAM J. Comput., 51(2):17–152, 2022. \href{https://doi.org/10.1137/17M1145288}{doi:10.1137/17M1145288}.

\bibitem[20]{bib20} John Canny. Some algebraic and geometric computations in PSPACE. In Proceedings of the Twentieth Annual ACM Symposium on Theory of Computing (STOC), pages 460–467, New York, NY, USA, 1988. ACM. \href{https://doi.org/10.1145/62212.62257}{doi:10.1145/62212.62257}.

\bibitem[21]{bib21} Thomas Colcombet, Nathanaël Fijalkow, Pawel Gawrychowski, and Pierre Ohlmann. The theory of universal graphs for infinite duration games. Log. Methods Comput. Sci., 18(3), 2022. \href{https://doi.org/10.46298/LMCS-18(3:29)2022}{doi:10.46298/LMCS-18(3:29)2022}.

\bibitem[22]{bib22} A. Condon. The complexity of stochastic games. Inform. and Comput., 96(2):203–224, 1992. \href{https://doi.org/10.1016/0890-5401(92)90048-K}{doi:10.1016/0890-5401(92)90048-K}.

\bibitem[23]{bib23} Anne Condon. On algorithms for simple stochastic games. In Jin-Yi Cai, editor, Advances in Computational Complexity Theory, volume 13 of DIMACS Series in Discrete Mathematics and Theoretical Computer Science, pages 51–72. AMS, 1993.

\bibitem[24]{bib24} Yann Disser, Oliver Friedmann, and Alexander V. Hopp. An exponential lower bound for Zadeh’s pivot rule. Math. Program., 199(1):865–936, 2023. \href{https://doi.org/10.1007/S10107-022-01848-X}{doi:10.1007/S10107-022-01848-X}.

\bibitem[25]{bib25} Andrzej Ehrenfeucht and Jan Mycielski. Positional strategies for mean payoff games. International Journal of Game Theory, 8(2):109–113, 1979.

\bibitem[26]{bib26} Friedrich Eisenbrand, Matthieu Haeberle, and Neta Singer. An Improved Bound on Sums of Square Roots via the Subspace Theorem. In Wolfgang Mulzer and Jeff M. Phillips, editors, 40th International Symposium on Computational Geometry (SoCG 2024), volume 293 of Leibniz International Proceedings in Informatics (LIPIcs), pages 54:1–54:8, Dagstuhl, Germany, 2024. Schloss Dagstuhl – Leibniz-Zentrum für Informatik. \href{https://doi.org/10.4230/LIPIcs.SoCG.2024.54}{doi:10.4230/LIPIcs.SoCG.2024.54}.

\bibitem[27]{bib27} E. A. Emerson, C. S. Jutla, and A. P. Sistla. On model-checking fragments of $\mu$-calculus. In Proceedings of the 5th International Conference on Computer-Aided Verification (CAV), volume 697 of Lecture Notes in Comput. Sci., pages 385–396. Springer, 1993. \href{https://doi.org/10.1007/3-540-56922-7_32}{doi:10.1007/3-540-56922-7\_32}.

\bibitem[28]{bib28} E. Allen Emerson and Charanjit S. Jutla. Tree automata, mu-calculus and determinacy. In Proceedings of the 32nd annual Symposium on Foundations of Computer Science (FOCS’91), pages 368–377, Washington, DC, USA, 1991. IEEE Computer Society.

\bibitem[29]{bib29} John Fearnley, Spencer Gordon, Ruta Mehta, and Rahul Savani. Unique end of potential line. Journal of Computer and System Sciences, 114:1–35, 2020. \href{https://doi.org/10.1016/J.JCSS.2020.05.007}{doi:10.1016/J.JCSS.2020.05.007}.

\bibitem[30]{bib30} John Fearnley, Marcin Jurdziński, and Rahul Savani. Linear complementarity algorithms for infinite games. In SOFSEM 2010: Theory and practice of computer science. 36th conference on current trends in theory and practice of computer science, Špindlerův Mlýn, Czech Republic, January 23–29, 2010. Proceedings, pages 382–393. Berlin: Springer, 2010. \href{https://doi.org/10.1007/978-3-642-11266-9_32}{doi:10.1007/978-3-642-11266-9\_32}.

\bibitem[31]{bib31} O. Friedmann. A subexponential lower bound for Zadeh’s pivoting rule for solving linear programs and games. In Proceedings of the 15th International Conference on Integer Programming and Combinatorial Optimization (IPCO), volume 6655 of Lecture Notes in Comput. Sci., pages 192–206. Springer, 2011. \href{https://doi.org/10.1007/978-3-642-20807-2_16}{doi:10.1007/978-3-642-20807-2\_16}.

\bibitem[32]{bib32} Oliver Friedmann, Thomas Dueholm Hansen, and Uri Zwick. Subexponential lower bounds for randomized pivoting rules for the simplex algorithm. In Proceedings of the Forty-third Annual ACM Symposium on Theory of Computing (STOC), pages 283–292, New York, NY, USA, 2011. ACM. \href{https://doi.org/10.1145/1993636.1993675}{doi:10.1145/1993636.1993675}.

\bibitem[33]{bib33} Oliver Friedmann, Thomas Dueholm Hansen, and Uri Zwick. Errata for: A subexponential lower bound for the random facet algorithm for parity games. CoRR, abs/1410.7871, 2014. \href{https://arxiv.org/abs/1410.7871}{arXiv:1410.7871}.

\bibitem[34]{bib34} M. R. Garey, R. L. Graham, and D. S. Johnson. Some NP-complete geometric problems. In Eighth Annual ACM Symposium on Theory of Computing (Hershey, Pa., 1976), pages 10–22. Assoc. Comput. Mach., New York, 1976.

\bibitem[35]{bib35} Bernd Gärtner and Leo Rüst. Simple stochastic games and P-matrix generalized linear complementarity problems. In Fundamentals of computation theory. 15th international symposium, FCT 2005, Lübeck, Germany, August 17–20, 2005. Proceedings., pages 209–220. Berlin: Springer, 2005. \href{https://doi.org/10.1007/11537311}{doi:10.1007/11537311}.

\bibitem[36]{bib36} M. Grötschel, L. Lovász, and L. Schrijver. Geometric Algorithms and Combinatorial Optimization. Springer, Heidelberg, 1994. Second edition.

\bibitem[37]{bib37} Thomas Dueholm Hansen and Rasmus Ibsen-Jensen. The complexity of interior point methods for solving discounted turn-based stochastic games. In The nature of computation. Logic, algorithms, applications. 9th conference on computability in Europe, CiE 2013, Milan, Italy, July 1–5, 2013. Proceedings, pages 252–262. Berlin: Springer, 2013. \href{https://doi.org/10.1007/978-3-642-39053-1_29}{doi:10.1007/978-3-642-39053-1\_29}.

\bibitem[38]{bib38} Wilfrid Hodges. A shorter model theory. Cambridge University Press, Cambridge, 1997.

\bibitem[39]{bib39} Marcin Jurdzinski, Rémi Morvan, and K. S. Thejaswini. Universal algorithms for parity games and nested fixpoints. In Jean-François Raskin, Krishnendu Chatterjee, Laurent Doyen, and Rupak Majumdar, editors, Principles of Systems Design - Essays Dedicated to Thomas A. Henzinger on the Occasion of His 60th Birthday, volume 13660 of Lecture Notes in Computer Science, pages 252–271. Springer, 2022. \href{https://doi.org/10.1007/978-3-031-22337-2_12}{doi:10.1007/978-3-031-22337-2\_12}.

\bibitem[40]{bib40} Marcin Jurdziński and Rahul Savani. A simple P-matrix linear complementarity problem for discounted games. In Logic and theory of algorithms. 4th conference on computability in Europe, CiE 2008, Athens, Greece, June 15–20, 2008. Proceedings, pages 283–293. Berlin: Springer, 2008. \href{https://doi.org/10.1007/978-3-540-69407-6_32}{doi:10.1007/978-3-540-69407-6\_32}.

\bibitem[41]{bib41} L. Khachiyan. A polynomial algorithm in linear programming. Doklady Akademii Nauk SSSR, 244:1093–1097, 1979.

\bibitem[42]{bib42} Jan Křetínský, Emanuel Ramneantu, Alexander Slivinskiy, and Maximilian Weininger. Comparison of algorithms for simple stochastic games. Information and Computation, 289:104885, 2022. \href{https://doi.org/10.1016/J.IC.2022.104885}{doi:10.1016/J.IC.2022.104885}.

\bibitem[43]{bib43} J. B. Lasserre. Moments, Positive Polynomials and Their Applications. Ser. Optim. Appl. Imperial College Press, 2009. \href{https://doi.org/10.1142/p665}{doi:10.1142/p665}.

\bibitem[44]{bib44} Karoliina Lehtinen and Udi Boker. Register games. Log. Methods Comput. Sci., 16(2), 2020. \href{https://doi.org/10.23638/LMCS-16(2:6)2020}{doi:10.23638/LMCS-16(2:6)2020}.

\bibitem[45]{bib45} Karoliina Lehtinen, Pawel Parys, Sven Schewe, and Dominik Wojtczak. A recursive approach to solving parity games in quasipolynomial time. Log. Methods Comput. Sci., 18(1), 2022. \href{https://doi.org/10.46298/LMCS-18(1:8)2022}{doi:10.46298/LMCS-18(1:8)2022}.

\bibitem[46]{bib46} T. M. Liggett and S. A. Lippman. Stochastic games with perfect information and time average payoff. SIAM Rev., 11(4):604–607, 1969. \href{https://doi.org/10.1137/1011093}{doi:10.1137/1011093}.

\bibitem[47]{bib47} Donald A. Martin. Borel determinacy. Annals of Mathematics, 102(2):363–371, 1975. URL: \href{http://www.jstor.org/stable/1971035}{http://www.jstor.org/stable/1971035}.

\bibitem[48]{bib48} Rolf H. Möhring, Martin Skutella, and Frederik Stork. Scheduling with and/or precedence constraints. SIAM Journal on Computing, 33(2):393–415, 2004. \href{https://doi.org/10.1137/S009753970037727X}{doi:10.1137/S009753970037727X}.

\bibitem[49]{bib49} A.W. Mostowski. Games with forbidden positions. Tech. Rep. 78, University of Gdansk, 1991.

\bibitem[50]{bib50} Y. Nesterov and A. Nemirovskii. Interior-point polynomial algorithms in convex programming, volume 13 of Studies in Applied and Numerical Mathematics. SIAM, Philadelphia, PA, 1994. \href{https://doi.org/10.1137/1.9781611970791}{doi:10.1137/1.9781611970791}.

\bibitem[51]{bib51} Lorant Porkolab and Leonid Khachiyan. On the complexity of semidefinite programs. Journal of Global Optimization, 10(4):351–365, 1997. \href{https://doi.org/10.1023/A:1008203903341}{doi:10.1023/A:1008203903341}.

\bibitem[52]{bib52} Anuj Puri. Theory of hybrid systems and discrete event systems. University of California at Berkeley, 1995.

\bibitem[53]{bib53} M. L. Puterman. Markov Decision Processes: Discrete Stochastic Dynamic Programming. Wiley Ser. Probab. Stat. Wiley, Hoboken, NJ, 2005.

\bibitem[54]{bib54} Motakuri V. Ramana. An exact duality theory for semidefinite programming and its complexity implications. Mathematical Programming, 77:129–162, 1997. \href{https://doi.org/10.1007/BF02614433}{doi:10.1007/BF02614433}.

\bibitem[55]{bib55} C. Scheiderer. Spectrahedral shadows. SIAM J. Appl. Algebra Geom., 2(1):26–44, 2018. \href{https://doi.org/10.1137/17M1118981}{doi:10.1137/17M1118981}.

\bibitem[56]{bib56} Sven Schewe. From parity and payoff games to linear programming. In Proceedings of the 34th International Symposium on Mathematical Foundations of Computer Science (MFCS), pages 675–686, 2009. \href{https://doi.org/10.1007/978-3-642-03816-7_57}{doi:10.1007/978-3-642-03816-7\_57}.

\bibitem[57]{bib57} M. Skomra. Tropical spectrahedra: Application to semidefinite programming and mean payoff games. PhD thesis, Université Paris-Saclay, 2018. URL: \href{https://pastel.archives-ouvertes.fr/tel-01958741}{https://pastel.archives-ouvertes.fr/tel-01958741}.

\bibitem[58]{bib58} Ola Svensson and Sergei Vorobyov. Linear complementarity and P-matrices for stochastic games. In Perspectives of Systems Informatics. 6th International Andrei Ershov Memorial Conference, PSI 2006, volume 4378 of Lecture Notes in Comput. Sci., pages 409–423. Springer, 2006. \href{https://doi.org/10.1007/978-3-540-70881-0_35}{doi:10.1007/978-3-540-70881-0\_35}.

\bibitem[59]{bib59} Sergey P. Tarasov and Mikhail N. Vyalyi. Semidefinite programming and arithmetic circuit evaluation. Discrete Applied Mathematics, 156(11):2070–2078, 2008. In Memory of Leonid Khachiyan (1952-2005). \href{https://doi.org/10.1016/J.DAM.2007.04.023}{doi:10.1016/J.DAM.2007.04.023}.

\bibitem[60]{bib60} Uri Zwick and Mike Paterson. The complexity of mean payoff games on graphs. Theor. Comput. Sci., 158(1\&2):343–359, 1996. \href{https://doi.org/10.1016/0304-3975(95)00188-3}{doi:10.1016/0304-3975(95)00188-3}.

\end{thebibliography}
\end{document}
